% Lösungen Prozentrechnung · Prozentwert · Prüfungs-Fokus MSA (zusammenbau v0.6)
\einheitenkopf{Einheit 3 · Prozentwert berechnen}

\erg{1}{a) $40\,€$ \quad b) $1\,\% = 4$ kg; $6\,\% = 24$ kg}
\erg{2}{a) $72\,€$ ($480 \cdot 0{,}15 = 72$; $408\,€$ ist der Restpreis) \quad b) $195\,€$ ($650 \cdot 0{,}3 = 195$; $455\,€$ ist der Restpreis) \quad c) $240 \cdot 0{,}15 = 36\,€$ (nicht der Restpreis $204\,€$) \quad d) $180 \cdot 0{,}35 = 63\,€$ (nicht der Restpreis $117\,€$) \quad e) $60 \cdot 0{,}17 = 10{,}20\,€$ \quad f) $40 \cdot 0{,}19 = 7{,}60\,€$ \quad g) $0{,}4 \cdot 85 = 34\,€$ (nicht $0{,}04$) \quad h) $0{,}6 \cdot 45 = 27\,€$ (nicht $0{,}06$) \quad i) $52\,\% - 35\,\% = 17$ Prozentpunkte; $0{,}17 \cdot 800 = 136$ Haushalte (nicht $17$) \quad j) $44\,\% - 26\,\% = 18$ Prozentpunkte; $0{,}18 \cdot 650 = 117$ Schüler (nicht $18$)}
